% !TeX root = ../泛函分析笔记.tex
% 由主文档通过 \input 载入；请编译主文档。
\setcounter{chapter}{2} % 保留原书章号。
\chapter{复分析与谱论的补充工具}
\label{app:fa-complex}

本附录只列出谱论直接使用的复分析工具，并说明算子值表达式的含义。复分析中的积分定理与逼近定理作为先修结论使用。

\section{实空间的复化}
若 $X$ 为实 Banach 空间，令 $X_{\mathbb C}=\{x+\mathrm iy:x,y\in X\}$，按通常规则定义复数乘法。一种可用范数是
\[
 \|x+\mathrm iy\|_{\mathbb C}=
 \sup_{\theta\in\mathbb R}\|x\cos\theta-y\sin\theta\|.
\]
旋转参数给出复齐次性，三角不等式逐个 $\theta$ 验证；并有
\[
 \max(\|x\|,\|y\|)\leq\|x+\mathrm iy\|_{\mathbb C}\leq\|x\|+\|y\|.
\]
所以它完备，且实空间等距嵌入。实有界算子 $T$ 延拓为 $T_{\mathbb C}(x+\mathrm iy)=Tx+\mathrm iTy$，其范数仍为 $\|T\|$。实算子的复谱定义为 $\sigma(T_{\mathbb C})$；实平面九十度旋转没有实特征值，却有复特征值 $\mathrm i,-\mathrm i$。

\section{解析函数的基本结论}
复函数在开集内每一点复可微时称解析（holomorphic）。本书使用 Cauchy 积分公式及其估计：若 $f$ 在包含闭圆盘 $|z-z_0|\leq R$ 的邻域内解析，则
\[
 f^{(n)}(z_0)=\frac{n!}{2\pi\mathrm i}
 \int_{|z-z_0|=R}\frac{f(z)}{(z-z_0)^{n+1}}\,\mathrm dz,
 \qquad |f^{(n)}(z_0)|\leq\frac{n!}{R^n}\max_{|z-z_0|=R}|f(z)|.
\]
若 $f$ 是有界整函数，让 $R\to\infty$ 即得 $f'=0$，所以 $f$ 为常数；这就是 Liouville 定理。若一条闭曲线在解析区域内连续变形，积分值不变。谱半径公式的证明使用了这一性质。

连续函数演算还使用 Weierstrass 逼近定理：实紧区间上的连续函数可由多项式一致逼近；对实紧子集先作连续延拓即可。复平面紧集上的一般连续函数则需使用含 $z,\overline z$ 的多项式及 Stone--Weierstrass 定理，不能不加条件地只用 $z$ 的多项式。

\section{算子值积分与标量化}
若 $F:[a,b]\to\mathcal B(X)$ 按算子范数连续，则其 Riemann 和在完备空间 $\mathcal B(X)$ 中收敛，定义积分，并满足
\[
 \left\|\int_a^bF(t)\,\mathrm dt\right\|\leq\int_a^b\|F(t)\|\,\mathrm dt.
\]
分段光滑曲线上的积分由参数化定义。每个 $x\in X$ 和 $f\in X^*$ 都可移入积分，即 $f((\int F)x)=\int f(Fx)$。因此可以先证明标量积分恒等式，再利用连续对偶分离向量恢复算子恒等式。

预解算子的 Neumann 展开按算子范数在足够小的闭圆盘上一致收敛，所以可逐项积分。在大圆上提取 $R(\lambda,T)=\sum_{k\geq0}T^k\lambda^{-k-1}$ 的第 $n$ 项，就得到谱半径证明中的积分表达。这里的一致收敛是交换求和与积分的依据。

\section{习题}
\begin{enumerate}
 \item 验证上述复化范数的复齐次性及完备性。
 \item 从预解恒等式推出 $\frac{\mathrm d}{\mathrm d\lambda}R(\lambda,T)=-R(\lambda,T)^2$。
 \item 在有限维对角矩阵上直接核验谱半径证明中的轮廓积分公式。
 \item 说明算子值函数弱意义上的逐点信息为何不足以无条件交换极限与算子范数。
\end{enumerate}

